% Propietario conservado: /Users/ruben/Documents/ChatGPT/jueces y controles/REVISION_ARGUMENTAL_APERTURAS_CIERRES_20260915/II/manuscript_es/sections/11_confluencia.tex
\section{Confluence de l’action, de l’aire, de l’information et de la gravitation}
\label{ii:sec:confluencia}

\subsection{Énergie et entropie de la cellule de demi-action}

Nous fixons une même section orientée d’action et ses sections
gravitationnelle et thermique correspondantes. La réciprocité précédente permet d’exprimer la même cellule au moyen de
l’action, de la durée et de la longueur. Avec
\(\ell_P^2=\hbar G/c^3\), \(t_P=\ell_P/c\) et
\(E_P=\hbar/t_P\), on a
\[
 \frac{c^4\ell_P}{G}=\frac{\hbar}{t_P}=E_P.
\]
L’égalité découle de la substitution de l’échelle aréale déjà construite. Dans la
réalisation d’horizon définie par \eqref{ii:eq:ii-horizonte-normalizado}, les quantités
se normalisent comme
\begin{equation}
 E_\Lambda(R)=\frac{c^4R}{2G},\qquad
 S_H(R)=\frac{k_B\pi R^2}{\ell_P^2},\qquad
 T_H(R)=\frac{\hbar c}{2\pi k_BR}.
 \label{ii:eq:ii-horizonte-normalizado}
\end{equation}
La température correspond à la normalisation d’horizon cosmologique
de cette réalisation. Le facteur \(2\pi\) et la définition du rayon
font partie de l’ensemble d’équations utilisé.

\begin{theorem}[Identité de la cellule gravitationnelle de demi-action]
Dans la réalisation \eqref{ii:eq:ii-horizonte-normalizado},
\begin{align}
 T_H(R)S_H(R)&=E_\Lambda(R),\\
 T_H(\ell_P)S_H(\ell_P)&=\frac{E_P}{2}
 =\frac{k_B\Theta_{\rm clk}\log3}{2},\\
 E_\Lambda(\ell_P)t_P&=\frac{\hbar}{2},\qquad
 E_\Lambda(\ell_P)T_{108}=\frac h2.
\end{align}
\end{theorem}
\begin{proof}
Multiplier la température par l’entropie annule la normalisation
thermique et donne \(\hbar cR/(2\ell_P^2)=c^4R/(2G)\).
En \(R=\ell_P\), l’expression est \(\hbar c/(2\ell_P)
=\hbar/(2t_P)=E_P/2\). L’identité thermique du cycle donne
le second membre. Enfin, \(T_{108}=2\pi t_P\) et
\(h=2\pi\hbar\) produisent les deux relations d’action.
\end{proof}

Le théorème réunit des quantités qui appartenaient auparavant à des descriptions
distinctes. L’énergie de la cellule, le produit de la température par l’entropie,
l’information tritique et l’action intégrée sur sa durée sont
des lectures compatibles sous la normalisation déclarée. La constante
gravitationnelle intervient dans l’échelle aréale ; l’action fixe l’énergie
de la période ; Boltzmann transporte l’information adimensionnelle vers la
droite d’entropie. La démonstration montre comment ces fonctions
s’articulent dans une même identité.

\subsection{Conversion entre énergie de décision et accroissement aréal}
\label{ii:sec:ii-conversion-informacion-area}

L’énergie par décision tritique et l’accroissement élémentaire d’aire
s’obtiennent au moyen de lecteurs distincts du développement précédent :
\begin{equation}
 E_P=k_B\Theta_{\rm clk}\log3=\frac{\hbar}{t_P},
 \qquad
 \delta\mathcal A=\gamma_{\HMT}a_0\ell_P^2.
 \label{ii:eq:ii-confluencia-cuanto-energia-area}
\end{equation}
On maintient la même section orientée et la branche positive
\(\hbar,G,\gamma_{\HMT},a_0,t_P>0\). La première quantité
appartient à la droite d’énergie et la seconde à la droite d’aire.
Leur quotient définit le coefficient de conversion
\begin{equation}
 \mathcal T_{I/A}
 :=\frac{E_P}{\delta\mathcal A}
 =\frac{k_B\Theta_{\rm clk}\log3}
        {\gamma_{\HMT}a_0\ell_P^2},
 \qquad [\mathcal T_{I/A}]=[\text{énergie}]/[\text{aire}].
 \label{ii:eq:ii-confluencia-tension}
\end{equation}
Cette définition explicite la normalisation nécessaire pour transporter
une énergie de décision vers une résolution géométrique. Tous ses facteurs
sont spécifiés avant de former le quotient.

\begin{proposition}[Expression gravitationnelle du coefficient de conversion]
\label{ii:prop:ii-confluencia-tension}
Dans la même réalisation dimensionnelle,
\begin{equation}
 \mathcal T_{I/A}
 =\frac{c^3}{\gamma_{\HMT}a_0G t_P}
 =\frac{c^4}{\gamma_{\HMT}a_0G\ell_P}.
 \label{ii:eq:ii-confluencia-tension-gravedad}
\end{equation}
\end{proposition}
\begin{proof}
Substituer \(E_P=\hbar/t_P\) et \(\ell_P^2=\hbar G/c^3\)
dans \eqref{ii:eq:ii-confluencia-tension} annule la section d’action
commune et donne le premier membre. L’identité \(\ell_P=ct_P\)
produit le second. On conserve la dépendance de \(t_P\) et
\(\ell_P\) à l’égard de la réalisation précédemment construite.
\end{proof}

Dans la coupe minimale de \eqref{ii:eq:ii-area-corte-informacion}, la
réalisation d’une décision uniforme par liaison publie
\begin{equation}
 S_{\rm corte}^{\rm fis}=81k_B\log3,
 \qquad
 E_{\rm corte}^{(1)}=81E_P
 =\Theta_{\rm clk}S_{\rm corte}^{\rm fis}.
 \label{ii:eq:ii-confluencia-corte-energia}
\end{equation}
La première égalité utilise la capacité démontrée et l’état produit
uniforme ; la seconde assigne l’énergie de décision du même cycle à
chacun de ses quatre-vingt-un liens. Multiplier l’identité thermique
par quatre-vingt-un prouve l’équation d’état de cette réalisation.
L’état sectoriel de trente-six liens conserve sa propre
fonction de partition, son entropie et les variations aréales de
\eqref{ii:eq:ii-area-ley-finita}.

La composition identifie les fonctions complémentaires des
constantes : l’action fixe l’énergie par durée ; le caractère angulaire
et Barbero--Immirzi fixent la résolution sectorielle d’aire ; Boltzmann
réalise dimensionnellement l’information ; le produit \(G\hbar\)
détermine l’unité aréale. Le coefficient \(\mathcal T_{I/A}\)
conserve ces dépendances en exprimant leur relation en unités
d’énergie par aire.

\subsection{Sens physique de la confluence aréale}

À l’horizon de rayon \(R\), l’aire est \(A_H=4\pi R^2\).
La même relation entropique admet les deux expressions
\[
 \frac{S_H}{k_B}=\frac{A_H}{4\ell_P^2}
 =\frac{\pi R^2c^3}{G\hbar}.
\]
Le premier quotient compare l’aire macroscopique à l’unité d’aire
de la construction. Le second explicite la composition de cette unité
par l’action et la gravitation. En transportant les sections de sorte
que \(G_a\hbar_a\) reste fixe, la lecture géométrique conserve
cette normalisation. L’entropie adimensionnelle du même horizon reste
invariante ; ses publications dimensionnelles conservent les bases qu’elles
utilisent.

Le paramètre de Barbero--Immirzi agit à une résolution différente
de ce même développement : l’incidence orientée organise le spectre
élémentaire d’aire et l’état réduit décrit l’information
accessible dans la coupure. Les feuillets aréal et volumétrique apportent le
transport géométrique et la mesure des histoires décrits dans
\cref{ii:sec:ii-hojas-av}. La comparaison entre un spectre microscopique
et une entropie d’horizon conserve comme données son application de réalisation,
l’état et l’échelle. Ces précisions permettent d’interpréter la
relation physique à chaque niveau et évitent de remplacer une correspondance
d’observables par une identification indifférenciée de leurs espaces.

L’unité théorique de l’argument réside dans des fonctions complémentaires :

L’unité d’aire réunit le produit $G\hbar$ ; l’incidence orientée
organise sa résolution élémentaire au moyen de $\gamma_{\HMT}a_0$ ;
l’état et la bipartition déterminent l’information accessible au
sous-système. Les équations de réponse
\eqref{ii:eq:ii-respuesta-entropica}--\eqref{ii:eq:ii-respuesta-areal}
explicitent cette articulation dans la famille sectorielle. L’inversion
\eqref{ii:eq:ii-inversion-entropia} conserve l’entropie et transforme
l’aire selon~\eqref{ii:eq:ii-area-complementaria} : les deux lectures
distinguent des propriétés différentes de la même configuration.

Dans cette articulation,
la dynamique produit et conserve l’histoire ; le lecteur détermine quelle
information est accessible ; l’action relie l’évolution et l’énergie ;
Boltzmann réalise son expression thermique ; la gravitation intervient
dans l’échelle géométrique avec laquelle l’aire est mesurée. L’identité
\(T_HS_H=E_\Lambda\) montre une confluence exacte de ces fonctions
dans la réalisation déclarée. La théorie physique se formule ainsi à partir
des objets mathématiques et de leurs transformations effectives.

\subsection{Entropie spectrale et conservation sous transport}

L’entropie d’un état réduit dépend de son spectre et de la
bipartition qui définit la réduction. Pour préciser la continuation
vers la dualité, soient \(\mathcal H_A,\mathcal H_B\) les deux facteurs
et \(\rho_{AB}\) un état d’entropies finies. Un transport
factorisé \(U_A\otimes U_B\) produit
\[
 \rho'_A=U_A\rho_AU_A^*,\qquad S(\rho'_A)=S(\rho_A).
\]
La première identité découle de la définition de la trace partielle et la
seconde de l’invariance du spectre. Pour un transport unitaire
général, la formulation équivalente consiste à transporter également
l’inclusion de l’algèbre des observables du sous-système. Cette précision
identifie la donnée géométrique qui doit accompagner une affirmation
de conservation de l’intrication.

La réciprocité de l’ellipse et l’échange quantité de mouvement--enroulement
conservent les quantités indiquées dans leurs théorèmes respectifs. Le
transport de la bipartition détermine comment utiliser ces identités
dans une lecture d’entropie. De cette manière, la continuation modulaire
maintient le contenu informationnel de l’article et conserve séparés
le spectre énergétique, l’orientation de l’action et la réduction
algébrique qui définit l’intrication.

\subsection{Multiplicité effective et réalisation hyperbolique}

Dans la cellule précédente, \(S_H(\ell_P)/k_B=\pi\). La multiplicité
effective associée à cette entropie est
\[
 \Omega_{\rm eff}=\exp(S_H/k_B)=e^\pi.
\]
Si \(H\) est l’involution hyperbolique du TRIT, de valeurs propres
\(+1,-1\), l’exponentielle \(\exp(\pi H)\) a pour spectre
\(\{e^\pi,e^{-\pi}\}\). La coïncidence de la valeur propre expansive
avec la multiplicité effective situe une relation spectrale spécifique.
Son sens microscopique conserve comme donnée supplémentaire l’application entre
l’espace de l’horizon et la représentation hyperbolique. Cette
localisation permet de développer la correspondance sans confondre
la valeur spectrale avec la totalité de son espace de réalisation.
